\section{小流时间极限下的运动方程}
\label{sec:5}
让我们考虑下列小流时间表示：
\begin{align}
   &\delta_{\mu\nu}\left[\Bar{\psi}(x)\overleftrightarrow{\Slash{D}}\psi(x)
   +2m_0\Bar{\psi}(x)\psi(x)\right]
\notag\\
   &=d_2(t)\Tilde{\mathcal{O}}_{2\mu\nu}(t,x)
   +d_4(t)\Tilde{\mathcal{O}}_{4\mu\nu}(t,x)
   +d_5(t)\Tilde{\mathcal{O}}_{5\mu\nu}+O(t),
\label{eq:(5.1)}
\end{align}
以及
\begin{equation}
   \delta_{\mu\nu}m_0\Bar{\psi}(x)\psi(x)
   =e_5(t)\Tilde{\mathcal{O}}_{5\mu\nu}(t,x)+O(t),
\label{eq:(5.2)}
\end{equation}
其中已理解方程两边均减去了真空期望值。通过与导出式~\eqref{eq:(4.22)}
和~\eqref{eq:(4.23)} 完全相同的重正化群论证以及第~\ref{sec:4} 节的单圈计算，对~$t\to0$
我们有
\begin{align}
   d_2(t)&=-\frac{1}{(4\pi)^2}T(R)N_{\mathrm{f}}\frac{5}{3},
\label{eq:(5.3)}\\
   d_4(t)&=1+\frac{\Bar{g}(1/\sqrt{8t})^2}{(4\pi)^2}C_2(R)
   \left[-\frac{1}{2}+\ln(432)\right],
\label{eq:(5.4)}\\
   d_5(t)&=\frac{\Bar{m}(1/\sqrt{8t})}{m}
   2\left\{1+\frac{\Bar{g}(1/\sqrt{8t})^2}{(4\pi)^2}C_2(R)
   \left[3\ln\pi+2+\ln(432)\right]\right\},
\label{eq:(5.5)}
\end{align}
且 $e_5(t)=(1/2)d_5(t)$。我们注意到，由于式~\eqref{eq:(5.1)}
的左边正比于费米子场的运动方程，当位置~$x$ 不与位置空间中其他算符的位置重合时，可以把该组
合置为零（Schwinger--Dyson 方程）。这意味着我们可以在主公式~\eqref{eq:(4.70)}
中作替换（减去真空期望值已理解）
\begin{align}
   \Tilde{\mathcal{O}}_{4\mu\nu}(t,x)
   &\to\frac{1}{(4\pi)^2}T(R)N_{\mathrm{f}}\frac{5}{3}
   \Tilde{\mathcal{O}}_{2\mu\nu}(t,x)
\notag\\
   &\qquad{}
   -\frac{\Bar{m}(1/\sqrt{8t})}{m}2
   \left[1+\frac{\Bar{g}(1/\sqrt{8t})^2}{(4\pi)^2}C_2(R)
   \left(3\ln\pi+\frac{5}{2}\right)\right]\Tilde{\mathcal{O}}_{5\mu\nu}(t,x)
\label{eq:(5.6)}
\end{align}
［对 $\overline{\text{MS}}$ 方案须作代换~\eqref{eq:(4.77)}］，因为本文始终假定能动量张
量~$\left\{T_{\mu\nu}\right\}_R(x)$ 在关联函数中与其他算符分离。这使得能动量张量的表达式
更为简洁。

若关注能动量张量的迹的部分——即膨胀流的全散度——上述步骤给出
\begin{align}
   \delta_{\mu\nu}\left\{T_{\mu\nu}\right\}_R(x)
   &=\lim_{t\to0}
   \biggl(\frac{1}{2}b_0
   \left[
   G_{\rho\sigma}^a(t,x)G_{\rho\sigma}^a(t,x)
   -\left\langle G_{\rho\sigma}^a(t,x)G_{\rho\sigma}^a(t,x)\right\rangle
   \right]
\notag\\
   &\qquad\qquad{}
   -\left\{1+\frac{\Bar{g}(1/\sqrt{8t})^2}{(4\pi)^2}C_2(R)
   \left[3\ln\pi+8+\ln(432)\right]\right\}
\notag\\
   &\qquad\qquad\qquad\qquad{}
   \times\Bar{m}(1/\sqrt{8t})\left[
   \mathring{\Bar{\chi}}(t,x)\mathring{\chi}(t,x)
   -\left\langle
   \mathring{\Bar{\chi}}(t,x)\mathring{\chi}(t,x)\right\rangle
   \right]
   \biggr),
\label{eq:(5.7)}
\end{align}
它与迹反常~\eqref{eq:(2.11)} 十分类似。对无质量费米子，$\Bar{m}(1/\sqrt{8t})=0$，
能动量张量迹部分的表达式最终变得十分简单。
